\chapter{الأشكال المعيارية وأشكال مااس وصيغ التتبع}
\label{ch:modular-maass-trace-formulas}

\section*{نطاق الفصل وحالته}

ينقل هذا الفصل نظرية الأعداد التحليلية من المجاميع على الأعداد الصحيحة إلى
التحليل الطيفي على السطح الزائدي. الفكرة الموحدة هي أن تعامد عائلة من
الدوال الذاتية يحول متوسطات معاملات فورييه إلى حد قطري وإلى مجاميع
Kloosterman موزونة بنوى Bessel. نعالج أولًا الهندسة الزائدية والأشكال
المعيارية الهولومورفية وأشكال مااس، ثم نعرض صيغ Petersson وKuznetsov
بتطبيع مصرح به، ونختم بمدخل بنيوي إلى صيغة Selberg في النموذج المدمج.

\noindent\textbf{حالة الفصل:}
\textenglish{AUTHORED-DRAFT / NON-CITABLE}. فُتحت بوابة التأليف بعد
مراجعة مستقلة لحزمة الأدلة والتطبيعات، لكن المتن ونتائجه لا يصبحان قابلين
للاستشهاد قبل البناء والتدقيق والمراجعة اللاحقة واعتماد المالك.

\subsection*{نطاق التطبيع}

\begin{itemize}
  \item في صيغة Petersson نعمل مع \(SL_2(\mathbb Z)\)، وشخصية تافهة،
  ووزن زوجي \(k>2\).
  \item في صيغة Kuznetsov نعمل مع الزمرة المعيارية ووزن مااس صفر،
  ونبقي الطيف المستمر ظاهرًا.
  \item صيغة Selberg تعرض للنموذج المدمج فقط؛ الحالة غير المدمجة وحدود
  القطع والمبعثرات مؤجلة.
\end{itemize}

\subsection*{نواتج التعلم}

بنهاية الفصل يستطيع القارئ أن:
\begin{enumerate}
  \item يثبت ثبات القياس الزائدي تحت فعل \(SL_2(\mathbb R)\).
  \item يميز بين شرط التحويل المعياري وشرط الانعدام عند الرؤوس.
  \item يستخرج توسع فورييه لشكل معياري من دوريته.
  \item يعرّف حاصل Petersson ويشرح سبب ظهور العامل \(y^k\).
  \item يكتب مؤثرات Hecke عند المستوى \(1\) ويستنتج علاقة معاملات الشكل الذاتي.
  \item يميز الأشكال الهولومورفية عن أشكال مااس الذاتية للابلاس.
  \item يحدد الطرف الطيفي والطرف الحسابي في صيغ Petersson وKuznetsov.
  \item يشرح لماذا لا يجوز حذف طيف Eisenstein من Kuznetsov.
  \item يقرأ صيغة Selberg المدمجة بوصفها مساواة بين الطول والطيف.
\end{enumerate}

\section{الفضاء العلوي والهندسة الزائدية}

نكتب
\[
\mathbb H=\{z=x+iy\in\mathbb C:y>0\},
\qquad
d\mu(z)=\frac{dx\,dy}{y^2}.
\]
تؤثر المصفوفة
\[
\gamma=
\begin{pmatrix}a&b\\c&d\end{pmatrix}
\in SL_2(\mathbb R)
\]
بواسطة تحويل Möbius
\[
\gamma z=\frac{az+b}{cz+d}.
\]
ولا يتغير الفعل إذا استبدلنا \(\gamma\) بـ\(-\gamma\)، ولذلك فالمجموعة
الهندسية الفعلية هي \(PSL_2(\mathbb R)\).

\begin{definition}[عامل الأتمتة والقياس الزائدي]
\label{def:automorphy-factor-hyperbolic-measure}
\resultid{ANT-DEF-20-01}
\provedhere
نضع
\[
j(\gamma,z)=cz+d.
\]
ويُسمى
\[
ds^2=\frac{dx^2+dy^2}{y^2},
\qquad
d\mu(z)=\frac{dx\,dy}{y^2}
\]
المتر والقياس الزائديين على \(\mathbb H\).
\end{definition}

\begin{proposition}[هويات الفعل وثبات القياس]
\label{prop:hyperbolic-invariance}
\resultid{ANT-PROP-20-01}
\provedhere
لكل \(\gamma\in SL_2(\mathbb R)\) و\(z\in\mathbb H\):
\[
\Im(\gamma z)=\frac{\Im z}{|cz+d|^2},
\qquad
\frac{d(\gamma z)}{dz}=\frac1{(cz+d)^2}.
\]
ومن ثم
\[
d\mu(\gamma z)=d\mu(z).
\]
\end{proposition}

\begin{proof}
باستعمال \(ad-bc=1\) نحسب
\[
\gamma z-\overline{\gamma z}
=
\frac{(az+b)(c\bar z+d)-(a\bar z+b)(cz+d)}
{|cz+d|^2}
=\frac{z-\bar z}{|cz+d|^2},
\]
فتنتج هوية الجزء التخيلي. كما يعطي الاشتقاق
\[
\frac{a(cz+d)-c(az+b)}{(cz+d)^2}
=\frac{ad-bc}{(cz+d)^2}.
\]
إذن عامل جاكوبي الحقيقي هو \(|cz+d|^{-4}\)، بينما يتحول
\(y^2\) إلى \(y^2|cz+d|^{-4}\). بإلغاء العاملين نحصل على ثبات
\(dx\,dy/y^2\).
\end{proof}

\begin{remark}
هذا الثبات هو السبب البنيوي في أن حاصل الضرب \(L^2\) لأشكال مااس لا
يحتاج وزنًا إضافيًا، بينما يحتاج حاصل Petersson للأشكال من الوزن \(k\)
إلى العامل \(y^k\) كي يعوض مقدار عامل الأتمتة.
\end{remark}

\section{الأشكال المعيارية الهولومورفية}

لتكن
\[
\Gamma=SL_2(\mathbb Z).
\]
إنها مولدة بالتحويلين
\[
T:z\mapsto z+1,
\qquad
S:z\mapsto-\frac1z.
\]

\begin{definition}[الشكل المعياري وشكل الحدبة]
\label{def:holomorphic-modular-cusp-form}
\resultid{ANT-DEF-20-02}
\citedresult[\cite[Ch.~1]{DiamondShurman2005}]
ليكن \(k\) عددًا صحيحًا. الشكل المعياري الهولومورفي من الوزن \(k\)
للمجموعة \(\Gamma\) هو دالة هولومورفية \(f:\mathbb H\to\mathbb C\)
تحقق
\[
f(\gamma z)=j(\gamma,z)^k f(z)
\qquad(\gamma\in\Gamma),
\]
وتكون هولومورفية عند الرأس \(\infty\). وإذا كان حدها الثابت عند الرأس
صفرًا سميت \emph{شكل حدبة}. نرمز إلى الفضاءين بـ\(M_k(\Gamma)\)
و\(S_k(\Gamma)\).
\end{definition}

\begin{remark}[قيد الزوجية]
لأن \(-I\in\Gamma\) يعمل على \(\mathbb H\) عمل الهوية، فإن
\[
f(z)=f((-I)z)=(-1)^k f(z).
\]
لذلك \(M_k(\Gamma)=\{0\}\) عندما يكون \(k\) فرديًا. ولهذا يقتصر
نطاق صيغة Petersson في هذا الفصل على الوزن الزوجي.
\end{remark}

\begin{lemma}[التوسع عند الرأس]
\label{lem:modular-fourier-expansion}
\resultid{ANT-LEM-20-01}
\provedhere
إذا كان \(f\) هولومورفيًا على \(\mathbb H\) ويحقق \(f(z+1)=f(z)\)،
فإن له توسعًا
\[
f(z)=\sum_{n\in\mathbb Z}a_f(n)e(nz),
\qquad
e(w)=e^{2\pi i w},
\]
متقاربًا محليًا. وإذا كان \(f\) هولومورفيًا عند \(\infty\)، فإن
\[
f(z)=\sum_{n\ge0}a_f(n)e(nz),
\]
وهو شكل حدبة عند \(\infty\) إذا وفقط إذا \(a_f(0)=0\).
\end{lemma}

\begin{proof}
ضع \(q=e(z)=e^{2\pi iz}\). تحول الدورية الدالة \(f\) إلى دالة
هولومورفية \(F(q)\) على القرص المثقوب \(0<|q|<1\)، لأن اختيار
لوغاريتم آخر لـ\(q\) يغير \(z\) بعدد صحيح ولا يغير \(f(z)\).
متسلسلة Laurent للدالة \(F\) تعطي
\[
F(q)=\sum_{n\in\mathbb Z}a_f(n)q^n.
\]
الهولومورفية عند الرأس تعني أن الشذوذ عند \(q=0\) قابل للإزالة، فتختفي
القوى السالبة. والانعدام عند الرأس يكافئ \(F(0)=a_f(0)=0\).
\end{proof}

\begin{example}[متسلسلة دلتا]
الدالة
\[
\Delta(z)=q\prod_{n\ge1}(1-q^n)^{24}
=\sum_{n\ge1}\tau(n)q^n
\]
شكل حدبة من الوزن \(12\). لا نثبت معياريتها هنا؛ المثال يوضح أن معاملات
فورييه للأشكال المعيارية تحمل متتاليات حسابية عميقة.
\end{example}

\section{حاصل Petersson}

\begin{definition}[حاصل Petersson عند المستوى \(1\)]
\label{def:petersson-inner-product}
\resultid{ANT-DEF-20-04}
\citedresult[\cite[\S3]{KnightlyLi2006}]
إذا كان \(f,g\in S_k(\Gamma)\)، نضع
\[
\langle f,g\rangle
=
\int_{\Gamma\backslash\mathbb H}
f(z)\overline{g(z)}\,y^k\,d\mu(z).
\]
\end{definition}

\begin{proposition}[سلامة تعريف حاصل Petersson]
\label{prop:petersson-invariance}
\resultid{ANT-PROP-20-02}
\provedhere
الدالة
\[
f(z)\overline{g(z)}\,y^k\,d\mu(z)
\]
ثابتة تحت فعل \(\Gamma\)، ويتقارب التكامل عندما يكون \(f,g\) شكلي
حدبة.
\end{proposition}

\begin{proof}
لدينا
\[
f(\gamma z)\overline{g(\gamma z)}
=|cz+d|^{2k}f(z)\overline{g(z)}
\]
و
\[
\Im(\gamma z)^k=\frac{y^k}{|cz+d|^{2k}},
\]
بينما \(d\mu\) ثابت. فتثبت اللاتغيرية. أما عند الرأس، فإن غياب الحد
الثابت يعطي تناقصًا أسيًا من توسع فورييه؛ وهذا يغلب نمو عامل \(y^k\).
وعلى الجزء المدمج من المجال الأساسي يكون التكامل منتهيًا بالاستمرارية.
\end{proof}

\section{مؤثرات Hecke}

نعتمد عند المستوى \(1\) التطبيع
\[
(T_nf)(z)
=
n^{k-1}
\sum_{\substack{ad=n\\d>0}}d^{-k}
\sum_{b\bmod d}f\!\left(\frac{az+b}{d}\right).
\]
ويمكن كتابته بصورة أوضح:
\[
(T_nf)(z)
=n^{k-1}\sum_{\substack{ad=n\\d>0}}d^{-k}
\sum_{b=0}^{d-1}f\!\left(\frac{az+b}{d}\right).
\]

\begin{theorem}[علاقات Hecke الأساسية]
\label{thm:hecke-relations-level-one}
\resultid{ANT-THM-20-01}
\citedresult[\cite[\S5.5]{DiamondShurman2005}]
تحفظ المؤثرات \(T_n\) الفضاء \(S_k(\Gamma)\)، وهي ذاتية المرافق بالنسبة
إلى حاصل Petersson، وتحقق
\[
T_mT_n
=
\sum_{d\mid(m,n)}d^{k-1}T_{mn/d^2}.
\]
وبخاصة، إذا \((m,n)=1\)، فإن \(T_mT_n=T_{mn}\).
\end{theorem}

إذا كان
\[
f(z)=\sum_{r\ge1}a_f(r)e(rz),
\]
فإن حساب مجموع \(b\bmod d\) باستعمال التعامد المنتهي يعطي
\[
a_{T_nf}(r)
=
\sum_{d\mid(n,r)}d^{k-1}a_f\!\left(\frac{nr}{d^2}\right).
\]
فإذا كان \(f\) شكلًا ذاتيًا مطبعًا بـ\(a_f(1)=1\)، ينتج من أخذ
\(r=1\) أن القيمة الذاتية لـ\(T_n\) هي \(a_f(n)\)، ومن ثم
\[
a_f(m)a_f(n)
=
\sum_{d\mid(m,n)}d^{k-1}a_f\!\left(\frac{mn}{d^2}\right).
\]
هذه العلاقة هي البذرة الجبرية للجداء الأويلري، لكن بناء دوال \(L\)
الآلية وتحليلها مؤجل إلى الفصل الحادي والعشرين.

\section{أشكال مااس}

نستعمل مؤثر لابلاس الزائدي بإشارة غير سالبة:
\[
\Delta=-y^2\left(\frac{\partial^2}{\partial x^2}
+\frac{\partial^2}{\partial y^2}\right).
\]

\begin{definition}[شكل مااس الحدبي من الوزن صفر]
\label{def:maass-cusp-form}
\resultid{ANT-DEF-20-03}
\citedresult[\cite[Chs.~3--4]{Iwaniec2002}]
شكل مااس الحدبي للمجموعة \(\Gamma\) هو دالة ملساء
\(u:\mathbb H\to\mathbb C\) تحقق:
\begin{enumerate}
  \item \(u(\gamma z)=u(z)\) لكل \(\gamma\in\Gamma\).
  \item \(\Delta u=\lambda u\) لقيمة ذاتية \(\lambda\).
  \item \(u\in L^2(\Gamma\backslash\mathbb H,d\mu)\).
  \item حدها الثابت عند الرأس يساوي صفرًا.
\end{enumerate}
\end{definition}

\begin{proposition}[فصل المتغيرات وتوسع Fourier--Whittaker]
\label{prop:maass-fourier-whittaker}
\resultid{ANT-THM-20-02}
\citedresult[\cite[\S2]{Kuznetsov1981}]
إذا كان \(u\) شكل مااس حدبيًا ذا قيمة ذاتية
\[
\lambda=\frac14+t^2,
\]
فإن
\[
u(z)
=
\sqrt y\sum_{n\ne0}\rho_u(n)
K_{it}(2\pi|n|y)e(nx).
\]
\end{proposition}

\begin{remark}[لماذا يظهر \(K\)-Bessel]
الدورية في \(x\) تعطي توسعًا فورييريًا. عند إدخاله في معادلة
\(\Delta u=(1/4+t^2)u\)، تحقق معاملات \(y\) معادلة Bessel معدلة.
شرط \(L^2\) والانحدار عند الرأس يستبعد الحل المتنامي \(I_{it}\) ويبقي
\(K_{it}\). هذه حجة فصل متغيرات؛ أما التحليل الطيفي الكامل وتمام العائلة
فمدخل مقتبس.
\end{remark}

\begin{remark}[تطبيعان لا يجوز خلطهما]
يمكن تطبيع الشكل بـ\(\|u\|_2=1\)، أو تطبيع معاملات Hecke بحيث
\(\rho_u(1)=1\). لا يتطابق التطبيعان عمومًا. صيغة Kuznetsov تتغير
عواملها إذا انتقلنا بينهما، ولذلك سنصرح بالتطبيع في الصيغة نفسها.
\end{remark}

\section{مجاميع Kloosterman}

\begin{definition}[مجموع Kloosterman]
\label{def:kloosterman-sum}
\resultid{ANT-DEF-20-05}
\provedhere
للأعداد الصحيحة \(m,n\) والعدد \(c\ge1\)، نضع
\[
S(m,n;c)
=
\sum_{\substack{d\bmod c\\(d,c)=1}}
e\!\left(\frac{md+n\bar d}{c}\right),
\]
حيث \(d\bar d\equiv1\pmod c\).
\end{definition}

\begin{proposition}[تماثلات أولية]
\label{prop:kloosterman-elementary-symmetries}
\resultid{ANT-PROP-20-03}
\provedhere
لدينا
\[
S(m,n;c)=S(n,m;c),
\qquad
\overline{S(m,n;c)}=S(-m,-n;c).
\]
\end{proposition}

\begin{proof}
في الهوية الأولى نستبدل \(d\) بمعكوسه \(\bar d\)، وهو تبديل لمجموعة
البواقي المختزلة. وفي الثانية نأخذ مرافق كل حد ونستعمل
\(\overline{e(x)}=e(-x)\).
\end{proof}

\begin{theorem}[حد Weil]
\label{thm:weil-kloosterman}
\resultid{ANT-THM-20-03}
\citedresult[\cite[Cor.~11.12]{IwaniecKowalski2004}]
لدينا
\[
|S(m,n;c)|
\le
\tau(c)(m,n,c)^{1/2}c^{1/2}.
\]
\end{theorem}

هذا ادخار جذر تربيعي، حتى عامل القواسم والقاسم المشترك، عن الحد التافه
\(|S(m,n;c)|\le\varphi(c)\). لا يثبت الفصل حد Weil؛ عمقه جبري هندسي
ويتجاوز الأدوات المطورة قبله.

\section{صيغة Petersson}

لتكن \(\mathcal B_k\) قاعدة متعامدة لـ\(S_k(\Gamma)\)، واكتب
\[
f(z)=\sum_{n\ge1}a_f(n)e(nz).
\]

\begin{theorem}[صيغة Petersson للمستوى \(1\)]
\label{thm:petersson-trace}
\resultid{ANT-THM-20-04}
\citedresult[\cite[Cor.~3.12]{KnightlyLi2006}]
إذا كان \(k>2\) زوجيًا و\(m,n\ge1\)، فإن
\[
\frac{\Gamma(k-1)}{(4\pi\sqrt{mn})^{k-1}}
\sum_{f\in\mathcal B_k}
\frac{a_f(m)\overline{a_f(n)}}{\langle f,f\rangle}
=
\delta_{m,n}
+2\pi i^{-k}
\sum_{c\ge1}
\frac{S(m,n;c)}{c}
J_{k-1}\!\left(\frac{4\pi\sqrt{mn}}{c}\right).
\]
\end{theorem}

تساوي الصيغة بين تعامد طيفي لمعاملات فورييه وطرف حسابي. الحد
\(\delta_{m,n}\) هو القطر؛ وكل انحراف عنه تقيسه مجاميع Kloosterman
مع نواة \(J_{k-1}\). ليست الصيغة مجرد تطبيق لمتباينة تعامد منتهية:
النواة البسليّة تسجل هندسة المدارات غير القطرية.

\section{صيغة Kuznetsov}

صيغة Kuznetsov هي النظير غير الهولومورفي لصيغة Petersson، لكن يوجد فرق
حاسم: السطح \(\Gamma\backslash\mathbb H\) غير مدمج، ولذلك لا يكتمل
الطيف بأشكال مااس الحدبية وحدها؛ يجب إضافة طيف Eisenstein المستمر.

نستعمل الصيغة الحديثة المطابقة في سجل التطبيعات. لتكن \(\{u_j\}\)
قاعدة متعامدة معيارية من أشكال مااس الحدبية، ذات معاملات
\(\rho_j(n)\) ومعلمات \(t_j\). واكتب توسع Eisenstein غير الصفري
بمعاملات \(\tau(n,t)\). ولتكن \(h\) دالة اختبار زوجية هولومورفية في
شريط أعرض من \(|\Im t|\le1/2\)، ومتلاشية بما يكفي، وضع
\[
\mathcal J(x,t)
=
\frac{J_{2it}(x)-J_{-2it}(x)}{\sinh(\pi t)},
\qquad
H^+(x)
=
\frac{i}{\pi}\int_{\mathbb R}
\mathcal J(x,t)h(t)t\tanh(\pi t)\,dt.
\]

\begin{theorem}[Kuznetsov، النواة ذات الإشارة المتساوية]
\label{thm:kuznetsov-trace}
\resultid{ANT-THM-20-05}
\citedresult[\cite[Lemma~2.1]{KhanYoung2020}]
لـ\(m,n\ge1\)، وبالتطبيع المبين في جدول الفصل:
\[
\begin{aligned}
&\sum_j
\frac{\rho_j(n)\overline{\rho_j(m)}}{\cosh(\pi t_j)}h(t_j)
+\frac1{4\pi}\int_{\mathbb R}
\frac{\tau(n,t)\overline{\tau(m,t)}}{\cosh(\pi t)}
h(t)\,dt\\
&\qquad
=
\frac{\delta_{m,n}}{\pi^2}\int_{\mathbb R}
h(t)t\tanh(\pi t)\,dt
+\sum_{c\ge1}\frac{S(m,n;c)}{c}
H^+\!\left(\frac{4\pi\sqrt{mn}}{c}\right),
\end{aligned}
\]
\end{theorem}

\begin{remark}[تطبيع Eisenstein]
إذا عُرّف
\[
E(z,s)=\frac12\sum_{\Gamma_\infty\backslash\Gamma}
\Im(\gamma z)^s,
\]
فإن \(\tau(n,t)\) يتضمن عوامل غاما وزيتا ودالة القواسم. يمكن إعادة
كتابة الحد المستمر بدوال القواسم وعامل
\(|\zeta(1+2it)|^{-2}\)، لكن لا يجوز نسخه من مرجع ذي تطبيع مختلف.
\end{remark}

\begin{remark}[تحويل \(K\)]
التحويل \(H^+\) هنا هو تحويل \(J\) للنواة ذات الإشارة المتساوية.
تحويل \(K\)-Bessel يظهر في صيغة الإشارة المتعاكسة، ولا نخلط الصيغتين
لمجرد أن كليهما يحمل اسم Kuznetsov.
\end{remark}

\section{مدخل بنيوي إلى صيغة Selberg}

ننتقل هنا إلى سطح زائد مدمج
\[
M=\Gamma\backslash\mathbb H
\]
بلا رؤوس، ولا عناصر قطع مكافئ، ولا طيف مستمر. اكتب قيم لابلاس
\[
\lambda_j=\frac14+r_j^2.
\]
لتكن \(h\) زوجية، هولومورفية في شريط، ومتلاشية بالسرعة المطلوبة، ولنعتمد
زوج التحويل
\[
g(u)=\frac1{2\pi}\int_{\mathbb R}h(r)e^{-iru}\,dr,
\qquad
h(r)=\int_{\mathbb R}g(u)e^{iru}\,du.
\]

\begin{theorem}[صيغة Selberg: النموذج المدمج]
\label{thm:selberg-compact-trace}
\resultid{ANT-THM-20-06}
\citedresult[\cite[Thm.~4]{Marklof2004}]
إذا جرى الجمع على الجيوديسيات المغلقة البدائية \(\gamma\) ذات الأطوال
\(\ell_\gamma\)، فإن
\[
\sum_{j\ge0}h(r_j)
=
\frac{\operatorname{Area}(M)}{4\pi}
\int_{\mathbb R}h(r)\,r\tanh(\pi r)\,dr
+\sum_{\gamma\ {\rm primitive}}\sum_{n\ge1}
\frac{\ell_\gamma\,g(n\ell_\gamma)}
{2\sinh(n\ell_\gamma/2)}.
\]
\end{theorem}

\deferredresult{معالجة مستقلة للصيغة غير المدمجة وحدود المبعثرات}

الطرف الأيسر طيفي، والطرف الأيمن هندسي: حد الهوية يقاس بالمساحة،
والحدود غير التافهة تقاس بأطوال الجيوديسيات المغلقة وتكراراتها. هذه
المساواة هي النظير الزائدي لفكرة أن الأثر يمكن حسابه من القيم الذاتية
أو من نقاط التثبيت/أصناف الاقتران.

\begin{remark}[حدود النموذج]
لا يجوز عرض هذه الصيغة بوصفها الصيغة الكاملة لـ
\(PSL_2(\mathbb Z)\backslash\mathbb H\)، لأن هذا السطح غير مدمج وله
رأس وطيف مستمر ومؤثر مبعثرات. عرضنا هنا نموذجًا مدمجًا مقصودًا لتثبيت
البنية من دون إخفاء الحدود الإضافية.
\end{remark}

\section{مقارنة الصيغ الثلاث}

\begin{center}
\begin{tabular}{llll}
\textbf{الصيغة} & \textbf{العائلة الطيفية} & \textbf{الطرف الحسابي/الهندسي}
& \textbf{النواة}\\
Petersson & أشكال هولومورفية & Kloosterman & \(J_{k-1}\)\\
Kuznetsov & مااس + Eisenstein & Kloosterman & تحويل \(J\)\\
Selberg المدمجة & طيف لابلاس & جيوديسيات مغلقة & \(g(n\ell)\)
\end{tabular}
\end{center}

التشابه هو مبدأ الأثر، لا تطابق الصيغ. صيغة Petersson مجموع محدود في
وزن ثابت، وKuznetsov تحليل طيفي غير مدمج، وSelberg تربط طيف لابلاس
بأصناف الاقتران الزائدية. نقل ثابت أو قياس من صف إلى آخر من دون قاموس
تطبيع يفسد النتيجة.

\section{كيف تخدم الصيغ نظرية الأعداد التحليلية؟}

\begin{enumerate}
  \item تحول متوسطات معاملات فورييه أو Hecke إلى حد قطري مفهوم ومجاميع
  Kloosterman يمكن تقديرها.
  \item تكشف الإلغاء الطيفي الذي لا يظهر من تقدير كل شكل على حدة.
  \item تدخل في لحظات دوال \(L\)، ومسائل دون التحدب، والغربال الطيفي،
  وتوزيع الأشكال الذاتية.
  \item تفصل بوضوح بين المعلومة المحلية عن معامل واحد والمعلومة
  المتوسطية عن عائلة كاملة.
\end{enumerate}

لكن هذه التطبيقات جسور إلى الفصل التالي، لا مدخلات لإثبات نتائج الفصل
الحالي. هذا يمنع اعتمادًا دائريًا على دوال \(L\) الآلية أو لانجلاندز.

\section{أخطاء شائعة}

\begin{enumerate}
  \item حذف عامل \(y^k\) من حاصل Petersson أو استعماله في حاصل مااس.
  \item الخلط بين \(a_f(n)\) و\(\lambda_f(n)\) من دون عامل التطبيع.
  \item الخلط بين التطبيع \(L^2\) وتطبيع \(\rho_j(1)=1\).
  \item حذف طيف Eisenstein من Kuznetsov للزمرة المعيارية.
  \item استعمال تحويل \(K\) في نواة الإشارة المتساوية.
  \item نقل عامل \(2\pi\) في تحويل Selberg إلى اصطلاح فورييري آخر.
  \item عرض النموذج المدمج لـSelberg بوصفه الصيغة غير المدمجة الكاملة.
  \item اعتبار حد Weil نتيجة مثبتة هنا لمجرد سهولة صياغتها.
\end{enumerate}

\section{تمارين}

\subsection{المستوى الأول}
\begin{enumerate}
  \item تحقق مباشرة من \(\Im(-1/z)=\Im(z)/|z|^2\).
  \item أثبت أن \(M_k(\Gamma)=\{0\}\) للأوزان الفردية.
  \item استخرج توسع \(q\) من الدورية \(f(z+1)=f(z)\).
  \item أثبت تماثلي مجموع Kloosterman المسجلين في الفصل.
\end{enumerate}

\subsection{المستوى الثاني}
\begin{enumerate}
  \item تحقق من ثبات تكامل Petersson تحت تغيير المجال الأساسي.
  \item احسب معاملات \(T_pf\) عندما يكون \(p\) أوليًا.
  \item اشتق معادلة Bessel المعدلة للمعامل الفورييري غير الصفري لشكل مااس.
  \item فسر لماذا يؤدي انعدام الحد الثابت إلى إزالة الجزء غير الحدبي.
\end{enumerate}

\subsection{المستوى الثالث}
\begin{enumerate}
  \item قارن تطبيع صيغة Petersson هنا بتطبيع يستعمل معاملات Hecke
  \(\lambda_f(n)\)، واكتب قاموس التحويل.
  \item حدد الحد الذي يفقد إذا حذفت طيف Eisenstein من Kuznetsov.
  \item ارسم خريطة تبين كيف ينتقل متوسط طيفي إلى حد قطري ومجموعات
  Kloosterman.
  \item اشرح الحدود الإضافية المتوقعة عند الانتقال من نموذج Selberg
  المدمج إلى السطح المعياري غير المدمج، من دون اشتقاقها.
\end{enumerate}

\section*{حالة نتائج المسودة}

\begin{center}
\begin{tabular}{ll}
\texttt{ANT-DEF-20-01} & \texttt{AUTHORED-DRAFT / PROVED-HERE}\\
\texttt{ANT-DEF-20-02} & \texttt{AUTHORED-DRAFT / CITED-DEFINITION}\\
\texttt{ANT-LEM-20-01} & \texttt{AUTHORED-DRAFT / PROVED-HERE}\\
\texttt{ANT-THM-20-01} & \texttt{AUTHORED-DRAFT / CITED}\\
\texttt{ANT-DEF-20-03} & \texttt{AUTHORED-DRAFT / CITED-DEFINITION}\\
\texttt{ANT-THM-20-02} & \texttt{AUTHORED-DRAFT / CITED}\\
\texttt{ANT-THM-20-03} & \texttt{AUTHORED-DRAFT / CITED}\\
\texttt{ANT-THM-20-04} & \texttt{AUTHORED-DRAFT / CITED-CORE}\\
\texttt{ANT-THM-20-05} & \texttt{AUTHORED-DRAFT / CITED-CORE}\\
\texttt{ANT-THM-20-06} & \texttt{AUTHORED-DRAFT / COMPACT-PROTOTYPE}
\end{tabular}
\end{center}

تبقى جميع النتائج \textenglish{NON-CITABLE} حتى اكتمال بناء PDF
والتدقيق الرياضي والمرجعي والمراجعة المستقلة اللاحقة واعتماد المالك.
